% !TeX root = ../main.tex
\begin{conclusions}
围绕传统粒子滤波在弱过程噪声下易塌缩为单点、在高维观测下似然权重失效这两类问题，本文结合最优输运与神经网络流形表示开展了理论分析与算法设计，主要结论如下：
\begin{enumerate}[label=(\arabic*),leftmargin=0pt,labelindent=2\ccwd,itemindent=*,labelsep=0.4em,nosep]
  \item 在低维离散重采样层面，建立了粒子重采样更新（涵盖 BPF 与 ETPF）的四项均方误差正交分解，并证明零过程噪声与确定性演化下自举粒子滤波（BPF）的不同粒子位置数期望按指数速率衰减至 $1$；仿真表明，强过程噪声下 BPF 即可低开销维持精度，而零过程噪声混沌系统中 BPF 全部塌缩为单点，ETPF 则完全避免了单点退化并显著降低估计误差。
  \item 在高维观测降维层面，提出了经验模内积下的压缩与还原流形拟合方法，给出平直线性子空间投影的伴随闭式解与谱截断性质，并推广至神经网络参数化的低维子流形，拟合高维观测对象流形以提取紧凑潜坐标并导出数据驱动的潜空间先验。
  \item 在潜空间后验更新层面，引入无需计算似然权重的条件输运映射对抗学习，利用潜状态静止时滤波后验的一步累积性质，提出离线统一训练、在线单次前向推理的静态摊销条件输运滤波算法；在含噪手写数字逐窗口识别实验中，传统加权滤波因权重集中而准确率偏低，逐步在线训练网络虽能提升精度但单步耗时较长，所提静态摊销方法同时实现了接近全图上限的识别准确率与亚毫秒级的单步推理速度。
\end{enumerate}

本文研究仍存在一定局限，后续工作可从以下方面继续推进：
\begin{enumerate}[label=(\arabic*),leftmargin=0pt,labelindent=2\ccwd,itemindent=*,labelsep=0.4em,nosep]
  \item 静态摊销方法的一步更新建立在潜状态不随时间变化的假设之上；面向状态随时间演化的动态系统，后续需引入循环网络或集合编码器汇总历史观测以构造时变条件输运映射。
  \item 离线摊销训练利用了训练集的成对样本，与仅在线访问有限先验粒子的传统加权滤波在已知信息量上并不对等，少样本或测试分布偏移条件下的映射泛化性能仍有待探索。
  \item 高维图像实验缺乏解析真值后验，目前主要借助分类误差与像素还原误差作间接评估，且在线对抗输运基线受限于固定训练步数预算；未来可结合可求解析后验的中高维基准及更大规模的重复实验作进一步检验。
\end{enumerate}
\end{conclusions}
